\documentclass[10pt,a4paper]{amsart}
\usepackage[margin=19mm]{geometry}
\usepackage{amsmath,amssymb,mathtools}
\usepackage[scr=rsfs,cal=euler]{mathalfa}
\usepackage{xfrac}
\usepackage[unicode,hidelinks,bookmarks=false]{hyperref}
\newcommand{\ie}{i.e.\ }
\newcommand{\cf}{cf.\ }
\newcommand{\resp}{respectively\ }
\newcommand{\Subsection}{\subsection}
\providecommand{\nobreakdash}{\nobreak\hbox{-}}
\setlength{\emergencystretch}{2em}
\multlinegap0pt
\usepackage{amssymb}
\usepackage{xcolor}
\newtheorem{lemma}{Lemma}
\def\a{{\alpha}} \def\bfalp{{\boldsymbol \alpha}} \def\atil{{\widetilde \alpha}}
\def\d{{\,{\rm d}}}
\def\eps{\varepsilon}
\def\grm{{\mathfrak m}}\def\grM{{\mathfrak M}}
\title{On weighted forms in many variables}
\author{Daniel Flores Galiote, Kiseok Yeon}
\date{Source: arXiv:2608.28774v1 (CC BY 4.0)}
\begin{document}
\maketitle

\noindent\emph{Source and licence.} On weighted forms in many variables. Version arXiv:2608.28774v1 (CC BY 4.0), distributed by arXiv under the Creative Commons Attribution 4.0 International licence, http://creativecommons.org/licenses/by/4.0/, as recorded in the arXiv licence metadata for this exact version and shown on the abstract page https://arxiv.org/abs/2608.28774v1. Full licence notice: \texttt{licenses/arxiv-2608.28774-CC-BY-4.0.txt}.\par\smallskip

\noindent\textit{Editorial scope.} Counted unit: the article's lemma bound (source0.tex lines 2094--2101). The article's complete original proof of it is delivered with it (source0.tex lines 2102--2170, one \texttt{proof} environment of 69 lines). No mathematics is written, repaired or re-derived: the statement and the proof are the article's own lines, in the article's own notation, and no retained line is substituted or re-flowed.  Retained uncounted as the source-local context the proof needs: lines 1148--1148; lines 2035--2064; lines 2058--2061.  Nothing was omitted from any retained span, so the package carries no omission record.  No retained line contains a citation, so the wrapper carries no bibliography and no bibliography entry.  No retained line contains a figure environment, an image-inclusion command or drawing code, so no image is delivered and the package declares no asset dependency.  Editorial formatting only, in addition: an AMS \texttt{amsart} preamble that restates the article's own declarations and macros used by the retained lines, namely the environments \texttt{lemma} on separate counters, together with the standard packages those lines need.  The retained environments print in this document as Lemma 1 in order of appearance, whereas the article prints the same items with the numbering of its own full document.  The complete native source of the article as distributed in the arXiv e-print for this exact version is bundled unchanged as \texttt{sources/208807/source0.tex}.
\section{The circle method}\label{sec:circle method}

\begin{lemma}\label{lemma:hat sigma pointwise minor arc bound}
    Suppose that
    \[
        X=P^{w_1/d},
        \qquad
        Y=P^{w_2/d}.
    \]
    Put
    \[
        \eta_*
        :=\min_{0\leq k\leq m}
        \begin{cases}
        \dfrac{w_2}{d}\bigl(D(k)+1\bigr),&d_1(k)=0,\\[3mm]
        \dfrac{w_1}{d}\bigl(D(k)+1\bigr),&d_1(k)>0,
        \end{cases}
    \]
    and set
    \[
        K:=K(F)=\max_{0\leq k\leq m}
        \frac{s_1+s_2-\widehat{\sigma}(k)}{\varkappa(k)}.
    \]
    If $K>0$, then, for every $0<\eta\leq\eta_*$, every $\eps>0$, and
    every $\a$ in the minor arcs $\grm(\eta)$, one has
    \begin{equation}\label{eqs: pointwise minor arc bound}
        |S(\a)|
        \ll X^{s_1}Y^{s_2}P^{-K\eta+\eps},
    \end{equation}
    where $\grm(\eta)=[0,1)\setminus\grM(\eta)$ and $\grM(\eta)$ is the
    set of major arcs defined in Section \ref{sec:circle method}.
\end{lemma}

    \begin{equation}
        |S(\a)|
        \ll X^{s_1}Y^{s_2}P^{-K\eta+\eps},
    \end{equation}

\begin{lemma}\label{lemma: minor arc bound}
    Suppose that $K > \max\{ 2, 1/\eta_*\}$. Then there exist $\delta>0$ and
    \[
        0<\vartheta_0<\frac{w_1}{5d},
    \]
    such that
    \[\int_{\a \in \grm(\vartheta_0)} |S(\a)| \d \a \ll X^{s_1}Y^{s_2}P^{-1-\delta}.\]
\end{lemma}

\begin{proof}
    Choose $\vartheta_0>0$ sufficiently small that
    \[
        \vartheta_0<\min\left\{\frac{w_1}{5d},\frac1K\right\},
        \qquad
        (K-2)\vartheta_0<K\eta_*-1,
    \]
    and then choose $\delta>0$ so that
    \[
        2\delta<(K-2)\vartheta_0.
    \]
    Put
    \[
        a=\left(1-\frac{2}{K}\right)\vartheta_0
        -\frac{2\delta}{K}>0,
        \qquad
        \vartheta_t=\vartheta_0+at.
    \]
    Let $T$ be the least positive integer for which
    \[
        K\vartheta_T\geq1+2\delta.
    \]
    The minimality of $T$ gives
    \[
        \vartheta_T<\frac{1+2\delta}{K}+a
        =\frac1K+\left(1-\frac2K\right)\vartheta_0
        <\eta_*,
    \]
    where the final inequality follows from the choice of $\vartheta_0$.
    Thus, the result of Lemma
    \ref{lemma:hat sigma pointwise minor arc bound} is available at every
    $\vartheta_t$ with $0\leq t\leq T$.

    By the definition of the major arcs, we have
    \[
        \grm(\vartheta_0)
        =\grm(\vartheta_T)\sqcup
        \bigsqcup_{1\leq t\leq T}
        \left(\grM(\vartheta_t)\setminus
        \grM(\vartheta_{t-1})\right).
    \]
    Moreover,
    \[
        \operatorname{meas}(\grM(\vartheta))
        \ll P^{2\vartheta-1}.
    \]
    The set
    $\grM(\vartheta_t)\setminus\grM(\vartheta_{t-1})$ is contained in
    $\grm(\vartheta_{t-1})$. Hence, by
    \eqref{eqs: pointwise minor arc bound},
    \begin{align*}
        \int_{\grm(\vartheta_0)}|S(\a)|\,\d\a
        &\ll X^{s_1}Y^{s_2}P^\eps
        \left(
            P^{-K\vartheta_T}
            +P^{-1}\sum_{1\leq t\leq T}
            P^{2\vartheta_t-K\vartheta_{t-1}}
        \right).
    \end{align*}
    Since $K\vartheta_T\geq1+2\delta$, the first term in parentheses is
    at most $P^{-1-2\delta}$. Furthermore, $K>2$ and
    \[
        2\vartheta_t-K\vartheta_{t-1}
        =(2-K)\vartheta_{t-1}+2a
        \leq(2-K)\vartheta_0+Ka
        =-2\delta.
    \]
    Since $T\ll1$, relabelling $\eps$ now gives the required estimate.
\end{proof}

\end{document}
